\section{Consequences of recent preprints}\label{sec:preprints}

Seven preprints released by OpenAI in September and October 2026
\cite{OAI26a,OAI26b,OAI26c,OAI26d,OAI26e,OAI26f,OAI26g} announce cases of the
Hodge conjecture that bear on (F2) and \textup{(F3$'$)}. They have not been
refereed, and we do not check their proofs here. Every statement of this
section that cites one of them is conditional on the theorem cited, and no
result outside this section depends on it. What we add are the deductions,
with complete proofs: a product argument that lowers the dimension of a Weil
class (\Cref{prop:stabilise}); the Hodge conjecture on every product of
powers of a K3 surface and of its Kuga--Satake variety, through a spanning
lemma for Clifford algebras (\Cref{lem:cliffordspan}, \Cref{thm:kspowers});
the transfer to fourfolds of K3$^{[2]}$ type whose transcendental part is
that of a K3 surface, which happens exactly when the N\'eron--Severi space
is isotropic or represents $-2$ (\Cref{prop:k3twosmall}); and the effect of the
preprints on the statements of the earlier sections. The theorems we use are
the following.
\begin{enumerate}[label=\textup{(P\arabic*)},leftmargin=2.8em,itemsep=2pt]
\item \cite[Theorem~1.1]{OAI26e}: the Hodge conjecture holds in every
  codimension on every complex abelian variety of CM type.
\item \cite[Theorem~1.1]{OAI26d}: for every imaginary quadratic field $K$ the
  Weil classes are algebraic on every abelian eightfold of
  $(K,-1,4)$-Weil type whose hermitian form has a totally isotropic
  $K$-subspace of dimension four.
\item \cite[Theorem~1.1]{OAI26c}: for every projective K3 surface $S$ the
  Kuga--Satake class of $S$ is algebraic, in the sense of
  \Cref{ssec:mumford}, for every choice of the data of
  \Cref{def:ksvariety} and \Cref{prop:ksembed}.%
\footnote{In \cite{OAI26c} the map $\Psi$ of \Cref{prop:ksembed} is stated
with values in $H^{1}(A)\otimes H^{1}(A)\subset H^{2}(A\times A,\QQ)$,
$A=\mathrm{KS}(T(S)_{\QQ})$, through the identification of
$\End H^{1}(A,\QQ)$ with $H^{1}(A)\otimes H^{1}(A)(1)$ given by a
polarisation. For the class $\theta$ of the polarisation and $g=\dim A$, the
pairing $(x,y)\mapsto\int_{A}x\,y\,\theta^{g-1}$ on $H^{1}(A,\QQ)$ is a
nonzero multiple of the form dual to $\theta$, and cup product with
$\mathrm{pr}_{1}^{*}\theta^{g-1}$ carries $H^{1}\otimes H^{1}$ isomorphically
onto the K\"unneth component $H^{2g-1}\otimes H^{1}$ that represents
$\End H^{1}(A,\QQ)$; both maps and the inverse of the second are induced by
algebraic cycles \cite[\S2]{Kle68}. So the two formulations differ by an
algebraic correspondence and a nonzero factor.}
\item \cite[Theorem~1.1]{OAI26f}: the Hodge conjecture holds in every
  codimension on every product of projective K3 surfaces. For the powers of
  one surface this also follows from (P3) and the criterion
  \cite[Theorem~1.1]{OAI26b}.
\item \cite[Theorems~1.1 and~1.2]{OAI26a}: for a connected abelian cover of a
  curve, the Hodge conjecture holds on every power of the covering curve at
  every parameter of its full marked family outside a countable union of
  proper closed analytic tensor Hodge loci; and for complete intersections
  of $c\le2$ diagonal hypersurfaces of a common degree $d\ge2$, it holds on
  every power of the very general member of the smooth family, and for $c=1$
  on every power of every smooth member $\sum_{i}a_{i}x_{i}^{d}=0$, in every
  dimension.
\end{enumerate}
The preprint \cite{OAI26g} proves the Hodge conjecture on every power of an
abelian sixfold of split Weil type and of an abelian variety of dimension at
most five with an action of an imaginary quadratic field; we do not use it.

\subsection{Weil classes in dimensions six and eight}\label{ssec:weilsix}

The Weil class is multiplicative (\Cref{thm:weilmult}), and an abelian
surface of Weil type has its Weil classes in degree two, where they are
divisor classes. Multiplying by such a surface therefore moves a Weil class
up by one dimension, and an algebraic representative upstairs can be cut
back down. The argument is Schoen's \cite[Proposition~10]{Sch98}, by which
Markman passes from the split sixfolds to the fourfolds; we give it in the
coordinates of \Cref{ssec:products}.

\begin{proposition}[Lowering the dimension]\label{prop:stabilise}
Let $K=\QQ(\sqrt{-d})$, let $n\ge1$, let $A$ be an abelian $2n$-fold of
$(K,-1,n)$-Weil type with discriminant $\delta$, and let $B$ be an abelian
surface of $(K,-1,1)$-Weil type with discriminant $\delta$, for instance the
quaternionic model of \Cref{thm:quatweil} with $n=1$ and $b$ representing
$\delta$.
\begin{enumerate}[label=\textup{(\roman*)},leftmargin=2.6em,itemsep=2pt]
\item $A\times B$ is of $(K,-1,n+1)$-Weil type with trivial discriminant,
  and its hermitian form is an orthogonal sum of $n+1$ hyperbolic planes, so
  it has a totally isotropic $K$-subspace of dimension $n+1$.
\item If the Weil classes of $A\times B$ are algebraic, so are those of $A$.
\end{enumerate}
\end{proposition}

\begin{proof}
(i) By \Cref{lem:products}, $A\times B$ is of $(K,-1,n+1)$-Weil type and its
discriminant is $\delta\cdot\delta=\delta^{2}$, the class of a square, which
is a norm, so it is trivial. The hyperbolic plane, with matrix
$\left(\begin{smallmatrix}0&1\\1&0\end{smallmatrix}\right)$, has signature
$(1,1)$ and $(-1)^{1}\det=1$, so an orthogonal sum of $n+1$ of them has
signature $(n+1,n+1)$ and trivial discriminant; by the classification
recalled in \Cref{def:disc} the hermitian form of $A\times B$ is isometric
to it, and the first basis vectors of the $n+1$ planes span a totally
isotropic subspace of dimension $n+1$.

(ii) By \Cref{thm:weilmult},
$\omega(A\times B)=\mathrm{pr}_{A}^{*}\omega(A)\cup\mathrm{pr}_{B}^{*}\omega(B)$
for the $K$-valued Weil classes. The class $\omega_{1}(B)$ is nonzero, being
one of two generators of the Weil line of $B$, and it is a divisor class by
\Cref{lem:basecase}. By the Hodge index theorem the intersection form on
$\NS(B)_{\QQ}$ is nondegenerate, so there is a divisor class $\beta$ on $B$
with $c_{1}=\int_{B}\omega_{1}(B)\,\beta\ne0$; put
$c_{2}=\int_{B}\omega_{2}(B)\,\beta\in\QQ$. The map
$p(\xi)=\mathrm{pr}_{A*}\bigl(\xi\cup\mathrm{pr}_{B}^{*}\beta\bigr)$ from
$H^{*}(A\times B,\QQ)$ to $H^{*}(A,\QQ)$ is induced by an algebraic cycle,
the graph of $\mathrm{pr}_{A}$ intersected with the pull-back of a divisor
representing $\beta$, and by the projection formula
$p(\mathrm{pr}_{A}^{*}\alpha\cup\mathrm{pr}_{B}^{*}\gamma)=\alpha\int_{B}
\gamma\beta$. Applied to \eqref{eq:weilmult} with $A_{1}=A$ and $A_{2}=B$ it
gives
\[
  p\bigl(\omega_{1}(A\times B)\bigr)=c_{1}\,\omega_{1}(A)-d\,c_{2}\,\omega_{2}(A),
  \qquad
  p\bigl(\omega_{2}(A\times B)\bigr)=c_{2}\,\omega_{1}(A)+c_{1}\,\omega_{2}(A).
\]
The determinant of this system is $c_{1}^{2}+d\,c_{2}^{2}>0$, so
$\omega_{1}(A)$ and $\omega_{2}(A)$ are rational combinations of the two
left-hand sides, which are algebraic when the Weil classes of $A\times B$
are.
\end{proof}

\begin{theorem}[Weil classes in dimensions six and eight]\label{thm:weilsix}
Assume \textup{(P2)}, and let $K$ be an imaginary quadratic field.
\begin{enumerate}[label=\textup{(\roman*)},leftmargin=2.6em,itemsep=2pt]
\item The Weil classes are algebraic on every abelian sixfold of
  $(K,-1,3)$-Weil type, for every discriminant, and on every abelian
  eightfold of $(K,-1,4)$-Weil type with trivial discriminant.
\item For $n=3$ and every $\delta$, and for $n=4$ and $\delta$ trivial, the
  algebraic locus of \Cref{thm:main} is the whole domain,
  $\Sigma_{K}=\cD_{K}$, and the Hodge conjecture holds for the very general
  member of the family.
\item In \Cref{thm:vgvandermonde}, the Hodge conjecture holds on
  $X_{6,6}(\lambda)$, and on $X_{d,8}(\lambda)$ for $d\in\{3,4\}$, for every
  $N$ and every $\lambda$ outside a countable union of proper Zariski closed
  subsets of $U$.
\end{enumerate}
\end{theorem}

\begin{proof}
(i) An eightfold of trivial discriminant has a hermitian form that is a sum
of four hyperbolic planes, as in the proof of \Cref{prop:stabilise}(i), so
\textup{(P2)} applies to it. For a sixfold $A$ of discriminant $\delta$,
\Cref{prop:stabilise} with $n=3$ gives an eightfold $A\times B$ to which
\textup{(P2)} applies, and part (ii) of the proposition brings the Weil
classes down to $A$.

(ii) By (i) the Weil classes are algebraic at every member, and
\Cref{thm:main}(iv) gives the Hodge conjecture for the very general member.

(iii) In the proof of \Cref{thm:vgvandermonde}(iii) every Hodge class of
degree $r$ is a sum of images of powers of polarisations and of Weil classes
of abelian $r$-folds of Weil type for $\QQ(\zeta_{m})$, $m\in\{3,4,6\}$
(\Cref{thm:vgvandermonde}(i), (ii)); for $r=6$ these are algebraic by (i)
whatever the discriminant, and for $r=8$ and $d\in\{3,4\}$ the discriminant
is trivial by \Cref{thm:vgvandermonde}(ii), so (i) applies again.
\end{proof}

For $n=4$ and nontrivial discriminant, \Cref{prop:stabilise} reduces the
Weil classes to the split tenfolds, which no result available to us
reaches.

\subsection{Kuga--Satake varieties}\label{ssec:kspowers}

The second deduction is a statement about Clifford algebras. Keep the
notation of \Cref{def:clifford} for a nondegenerate quadratic space $(T,q)$
over $\QQ$ of dimension $t\ge1$, so that $v^{2}=q(v)$ in $C=C(T,q)$.

\begin{lemma}[Two-sided multiplications]\label{lem:cliffordspan}
The $\QQ$-linear endomorphisms of $C^{+}$ are spanned by the maps
$x\mapsto a\,x\,b$ with $a,b\in C$ homogeneous of the same parity.
\end{lemma}

\begin{proof}
Each such map preserves $C^{+}$ by \Cref{lem:cliffdim}, and it suffices to
prove the statement after extending scalars to $\bar{\QQ}$, since the maps
are defined over $\QQ$. Over $\bar{\QQ}$ choose an orthogonal basis
$f_{1},\dots,f_{t}$ of $T$ with $q(f_{i})=1$, and write $f_{R}$ for the
increasing products of \Cref{lem:cliffdim}. The elements $\pm f_{R}$ form a
group $\Gamma$ of order $2^{t+1}$, since $f_{i}^{2}=1$ and
$f_{i}f_{j}=-f_{j}f_{i}$ for $i\ne j$, and the pairs $(g,h)\in\Gamma^{2}$ of
the same parity form a group $\Gamma_{2}$ of order $2^{2t+1}$, acting on
$C^{+}_{\bar{\QQ}}$ by $x\mapsto gxh^{-1}$. Its character $\chi$ is
computed on the basis $f_{R}$, $|R|$ even. The element
$(\pm f_{S},\pm f_{U})$ sends $f_{R}$ to a multiple of
$f_{S\triangle R\triangle U}$, so it has trace zero unless $S=U$. For $S=U$,
moving each factor of $f_{S}$ past $f_{R}$ gives
$f_{S}f_{R}f_{S}^{-1}=(-1)^{|S||R|-|S\cap R|}f_{R}=(-1)^{|S\cap R|}f_{R}$ for
$|R|$ even, and
\[
  \sum_{|R|\ \mathrm{even}}(-1)^{|S\cap R|}
  =\tfrac12\sum_{R}\bigl(1+(-1)^{|R|}\bigr)(-1)^{|S\cap R|}
  =\tfrac12\Bigl(\sum_{R}(-1)^{|S\cap R|}+\sum_{R}(-1)^{|R\setminus S|}\Bigr),
\]
which is $2^{t-1}$ for $S=\emptyset$ and for $S=\{1,\dots,t\}$ and zero
otherwise. So $|\chi|^{2}=2^{2t-2}$ on the eight elements
$(\pm f_{S},\pm f_{S})$ with $S\in\{\emptyset,\{1,\dots,t\}\}$ and zero
elsewhere, and
$\langle\chi,\chi\rangle=8\cdot2^{2t-2}/2^{2t+1}=1$. The representation is
therefore irreducible, and by Burnside's theorem the operators of
$\Gamma_{2}$ span $\End(C^{+}_{\bar{\QQ}})$.
\end{proof}

For the rest of this subsection let $S$ be a projective K3 surface,
$T=T(S)_{\QQ}$ with the intersection form $q$, $t=\dim T$,
$W=C^{+}(T)=H^{1}(A,\QQ)$ for $A=\mathrm{KS}(T)$ as in \Cref{def:ksvariety},
$v_{0}\in T$ with $q(v_{0})\ne0$, and $\Psi$ the embedding of
\Cref{prop:ksembed}. Let $\pi_{1}$ be the K\"unneth projector of $A$ onto
$H^{1}$, which is algebraic \cite[\S2]{Kle68}. For a smooth projective variety
$X$ and a subspace $H\subseteq H^{*}(X,\QQ)$ we say that a linear map
$f\colon H\to\End_{\QQ}(W)$ is \emph{induced} by an algebraic cycle $Z$ on
$X\times A\times A$ if for every $y\in H$ the correspondence
$\pi_{1}\circ Z_{*}(y)\circ\pi_{1}$, where
$Z_{*}(y)=\mathrm{pr}_{23*}(Z\cdot\mathrm{pr}_{1}^{*}y)$, acts on $H^{1}(A,\QQ)$
as $f(y)$ and by zero on the other degrees. By (P3) and the footnote to it,
$\Psi$ is induced on $T\subseteq H^{2}(S,\QQ)$ by an algebraic cycle.

\begin{lemma}[The Kuga--Satake class generates]\label{lem:ksgenerates}
Assume that $\Psi$ is induced on $T$ by an algebraic cycle. For $j\ge0$ and
$b\in C^{+}$ let $\Phi_{j,b}\colon T^{\otimes j}\to\End_{\QQ}(W)$ be
\[
  \Phi_{j,b}(v_{1}\otimes\dots\otimes v_{j})
  =\Psi(v_{1})\circ\dots\circ\Psi(v_{j})\circ R_{b},
\]
where $R_{b}$ is right multiplication by $b$ and $T^{\otimes j}$ is the
K\"unneth component of $H^{2j}(S^{j},\QQ)$. Then $\Phi_{j,b}$ is induced by an
algebraic cycle on $S^{j}\times A\times A$, and the images of the
$\Phi_{j,b}$, $0\le j\le t$, $b\in C^{+}$, span $\End_{\QQ}(W)$.
\end{lemma}

\begin{proof}
The complex structure of $W$ is left multiplication by an element of
$C^{+}\otimes\RR$ (\Cref{lem:cliffJ}), which commutes with $R_{b}$; so $R_{b}$
is an endomorphism of the Hodge structure $W$, hence an element of
$\End^{0}(A)$, and a multiple of the graph of an endomorphism induces it.
If $f_{1}$ on $H_{1}\subseteq H^{*}(X_{1},\QQ)$ and $f_{2}$ on
$H_{2}\subseteq H^{*}(X_{2},\QQ)$ are induced by $Z_{1}$ and $Z_{2}$, then
$y_{1}\otimes y_{2}\mapsto f_{1}(y_{1})\circ f_{2}(y_{2})$ on
$H_{1}\otimes H_{2}\subseteq H^{*}(X_{1}\times X_{2},\QQ)$ is induced by the
cycle on $X_{1}\times X_{2}\times A\times A$ obtained from
$Z_{1}\times Z_{2}$ by intersecting with the pull-back of the diagonal of the
two middle factors $A$ and pushing forward, composed on both sides with
$\pi_{1}$: its value at $y_{1}\otimes y_{2}$ is the composite of the
correspondences $\pi_{1}\circ Z_{1*}(y_{1})\circ\pi_{1}$ and
$\pi_{1}\circ Z_{2*}(y_{2})\circ\pi_{1}$, which acts by the composite of their
actions \cite[\S16.1]{Ful98}. By induction each $\Phi_{j,b}$ is induced by
an algebraic cycle.

For the span, $\Psi(v)(x)=vxv_{0}$ and $v_{0}^{2}=q(v_{0})$ give
\[
  \Phi_{j,b}(v_{1}\otimes\dots\otimes v_{j})(x)=v_{1}\cdots v_{j}\,x\,b\,v_{0}^{j},
\]
where $v_{0}^{j}$ is a nonzero rational number for $j$ even and a nonzero
multiple of $v_{0}$ for $j$ odd. As $b$ runs over $C^{+}$, the element
$c=b\,v_{0}^{j}$ runs over $C^{+}$ for $j$ even and over $C^{+}v_{0}=C^{-}$
for $j$ odd, $v_{0}$ being invertible. For an orthogonal basis of $T$ the
products $v_{1}\cdots v_{j}$ include the basis elements $f_{R}$ with $|R|=j$
of \Cref{lem:cliffdim}, so as $j$ runs over the integers of one parity in
$[0,t]$ they span the part of $C$ of that parity. Hence the images contain
every map $x\mapsto axc$ with $a,c$ homogeneous of the same parity, and
\Cref{lem:cliffordspan} completes the proof.
\end{proof}

\begin{theorem}[Kuga--Satake varieties]\label{thm:kspowers}
Assume \textup{(P3)} and \textup{(P4)}, and let $S$ be a projective K3
surface with Kuga--Satake variety $A=\mathrm{KS}(T(S)_{\QQ})$. Then every Hodge
class on $A^{k}\times S^{l}$ is algebraic, for all $k,l\ge0$. The same holds
for $B^{k}\times S^{l}$ whenever $B$ is isogenous to an abelian subvariety of
a power of $A$.
\end{theorem}

\begin{proof}
\emph{Step 1.} Through Poincar\'e duality an endomorphism $f$ of $W$ is a
class $[f]$ in the K\"unneth component $H^{2g-1}(A)\otimes H^{1}(A)$ of
$H^{2g}(A\times A,\QQ)$, $g=\dim A$, and $[f]$ is the class of any
correspondence acting by $f$ on $H^{1}$ and by zero elsewhere. The inverse of
the Lefschetz isomorphism $\theta^{g-1}\colon H^{1}(A)\to H^{2g-1}(A)$ is
induced by an algebraic cycle \cite[\S2]{Kle68}, and applied to the first
factor it carries $H^{2g-1}(A)\otimes H^{1}(A)$ onto
$H^{1}(A)\otimes H^{1}(A)$. By \Cref{lem:ksgenerates}, $W\otimes W\subset
H^{2}(A\times A,\QQ)$ is therefore spanned by classes $Z_{*}(y)$ with
$y\in T^{\otimes j}$, $0\le j\le t$, and $Z$ an algebraic cycle on
$S^{j}\times A^{2}$.

\emph{Step 2.} Taking external products, $(W\otimes W)^{\otimes p}\subset
H^{2p}(A^{2p},\QQ)$ is spanned by classes $Z_{*}(y)$ with $y\in T^{\otimes J}$
and $Z$ algebraic on $S^{J}\times A^{2p}$. For $k\ge1$,
$H^{1}(A^{k},\QQ)=\bigoplus_{i}\mathrm{pr}_{i}^{*}W$ and
$H^{2p}(A^{k},\QQ)$ is spanned by products of $2p$ classes of degree one.
Grouping them in pairs, each product is $\mu^{*}$ of a class in
$(W\otimes W)^{\otimes p}$ for a morphism $\mu\colon A^{k}\to(A\times A)^{p}$
whose components are pairs of projections. So $H^{2p}(A^{k},\QQ)$ is spanned
by classes $Z_{*}(y)$ with $y\in T^{\otimes J}$ and $Z$ algebraic on
$S^{J}\times A^{k}$.

\emph{Step 3.} The cohomology of $S$ vanishes in odd degrees, so
$H^{2p}(A^{k}\times S^{l},\QQ)$ is the sum of the
$H^{2a}(A^{k})\otimes H^{2p-2a}(S^{l})$, and by Step 2 it is spanned by
finitely many subspaces $Z_{i*}(H_{i})$ with $H_{i}=T^{\otimes J_{i}}\otimes
H^{2p-2a_{i}}(S^{l})\subseteq H^{*}(S^{J_{i}+l},\QQ)$ and $Z_{i}$ algebraic.
Each $Z_{i*}$ is a morphism of Hodge structures up to a Tate twist, so their
sum is a surjective morphism of polarisable Hodge structures onto
$H^{2p}(A^{k}\times S^{l},\QQ)$, and since polarisable Hodge structures form
a semisimple category every Hodge class of degree $2p$ is
$\sum_{i}Z_{i*}(\eta_{i})$ with $\eta_{i}$ a Hodge class in $H_{i}$. Each
$\eta_{i}$ is a Hodge class on $S^{J_{i}+l}$, hence algebraic by
\textup{(P4)}, and so is its image. In odd degrees there are no Hodge
classes.

\emph{Step 4.} If $B$ is isogenous to an abelian subvariety of $A^{m}$, then
by Poincar\'e reducibility $A^{m}$ is isogenous to $B\times B'$ for some
$B'$. A Hodge class on $B^{k}\times S^{l}$ pulled back along the projection
from $(B\times B')^{k}\times S^{l}$ is a Hodge class, algebraic by Step 3
transported along the isogeny, and its restriction along
$b\mapsto(b,0)$ is the class itself.
\end{proof}

The hypothesis (P4) enters only through the powers of $S$; for a product of
several K3 surfaces with their Kuga--Satake varieties the argument would need
Hodge classes in $H^{1}(A_{1})\otimes H^{1}(A_{2})$ for two of the
varieties, which \Cref{lem:ksgenerates} does not reach.

\begin{corollary}[Abelian varieties of orthogonal type]\label{cor:kstype}
Assume \textup{(P3)} and \textup{(P4)}, and let $T$ be a polarised rational
Hodge structure of K3 type, as in \Cref{setup:k3type}, with $\dim T\le19$.
Then the Hodge conjecture holds on every power of $\mathrm{KS}(T)$ and of
every abelian variety isogenous to an abelian subvariety of a power of it.
\end{corollary}

\begin{proof}
Let $P\subseteq T$ be the smallest rational sub-Hodge structure whose
complexification contains $T^{2,0}$. It contains no nonzero rational class
of type $(1,1)$, since the orthogonal complement in $P$ of such a class
would be a smaller one, and $N=P^{\perp}\cap T$ consists of classes of type
$(1,1)$. As in \Cref{rem:mumfordpowersopen},
$C^{+}(T)=C^{+}(P)\otimes C^{+}(N)\oplus C^{-}(P)\otimes C^{-}(N)$ is, as a
Hodge structure, a sum of copies of $C^{+}(P)$, so $\mathrm{KS}(T)$ is
isogenous to a power of $\mathrm{KS}(P)$. The space $(P,q)$ has dimension at
most $19$ and signature $(2,\dim P-2)$, so by the argument of
\Cref{thm:mumfordks}(i) it embeds isometrically in $\Lambda_{\QQ}$, and the
surjectivity of the period map gives a projective K3 surface $S$ with
$T(S)_{\QQ}\cong P$ as Hodge structures with their forms. Then
$\mathrm{KS}(P)$ is the Kuga--Satake variety of $S$, and
\Cref{thm:kspowers} applies.
\end{proof}

\begin{corollary}[Mumford fourfolds]\label{cor:mumfordall}
Assume \textup{(P3)}, and let $X$ be an abelian fourfold of Mumford type, as
in \Cref{setup:mumford}. Then the Hodge conjecture holds on $X^{n}$ for every
$n\ge1$ and on $S_{\lambda}^{a}\times X^{b}$ for all $a,b\ge0$, with
$S_{\lambda}$ as in \Cref{thm:mumfordks}. In particular the two exceptional
classes of \Cref{prop:mumfordproduct}(ii) are algebraic, and for every family
$W\to C$ of \Cref{cor:mumfordlefschetz} the equivalent conditions of
\Cref{cor:mumfordequiv} and \Cref{cor:mumfordpowersequiv} hold.
\end{corollary}

\begin{proof}
By \textup{(P3)} the Kuga--Satake class of $S_{\lambda}$ is algebraic, and
\Cref{cor:mumfordksproducts} gives the first two statements. In a family
$W\to C$, a fibre $X_{t}$ that is not a CM point has Hodge group $G$, as
shown in the proof of \Cref{cor:mumfordequiv}, so it is as in
\Cref{setup:mumford} and the first statement applies to it; the CM fibres
are covered by \Cref{prop:mumfordcmpowers}. So condition (e) of
\Cref{cor:mumfordpowersequiv} holds.
\end{proof}

The Kuga--Satake variety of $S_{\lambda}$ is isogenous to $X^{32}$
(\Cref{thm:mumfordks}(i)), so \Cref{thm:kspowers} gives the same conclusion
from (P3) and (P4); \Cref{cor:mumfordksproducts} needs only (P3).

\subsection{K3 surfaces and \textup{(F3$'$)}}\label{ssec:k3all}

\begin{corollary}[K3 surfaces]\label{cor:k3all}
\leavevmode
\begin{enumerate}[label=\textup{(\roman*)},leftmargin=2.6em,itemsep=2pt]
\item Assume \textup{(P3)}. Every projective K3 surface lies in the class
  $\mathcal{A}$ of \Cref{prop:f3prime}(iii), and so do its Hilbert schemes of
  points, the varieties of \Cref{prop:aclosure}(v) built from K3 surfaces, and
  the blow-ups, projective bundles and products formed from these and the
  other members of $\mathcal{A}$; \textup{(F3$'$)} holds on all of them.
\item Assume \textup{(P4)}. The Hodge conjecture holds on every product of
  projective K3 surfaces and, for every projective K3 surface $S$, on every
  Hilbert scheme $S^{[n]}$ and every moduli space listed in
  \Cref{prop:k3powers}(i), whatever the field $\End_{\mathrm{Hdg}}(T(S))$.
\item Assume \textup{(P4)}. In every family of \Cref{setup:rm} the algebraic
  locus of the real multiplication is the whole domain,
  $\Sigma_{E}=\cD_{E}$, and the real multiplication of every projective K3
  surface is induced by algebraic cycles on $S\times S$. Of the frontier cases
  of \Cref{rem:f3primefrontier}, the varieties $S\times S$ and $S^{[2]}$ for a
  K3 surface with real multiplication, $S_{\lambda}$ among them, and the class
  $\det T(S)$ on $S^{21}$ for a very general K3 surface, with $S^{[231]}$, are
  settled.
\end{enumerate}
\end{corollary}

\begin{proof}
(i) The cycle of (P3) induces a nonzero map from $T(S)_{\QQ}$ to
$H^{2}(A\times A,\QQ)$, $A=\mathrm{KS}(T(S)_{\QQ})$, so \Cref{prop:aclosure}(iv)
puts $S$ in $\mathcal{A}$, and the rest of the statement is
\Cref{prop:aclosure} and \Cref{prop:f3prime}(iii).

(ii) The first statement is (P4). The proof of \Cref{prop:k3powers}(i) uses
the field $\End_{\mathrm{Hdg}}(T(S))$ only to obtain the Hodge conjecture on
the powers of $S$; from there the motivic decompositions
\cite{dCM02,Bue20,FFZ21} give the Hilbert schemes and the moduli spaces.

(iii) The class $\tilde e_{\omega}$ of \Cref{setup:rm} is a Hodge class on
$S_{\omega}\times S_{\omega}$, algebraic by (ii) at every $\omega$; and the
endomorphisms of $T(S)_{\QQ}$ given by the field act through Hodge classes on
$S\times S$. The class $\det T(S)$ is a Hodge class on $S^{21}$, and
$S^{[231]}$ is covered by (ii).
\end{proof}

The same hypothesis reaches the fourfolds of K3$^{[2]}$ type whose
transcendental part is that of a K3 surface, through Markman's theorem on
rational Hodge isometries.

\begin{proposition}[Fourfolds of K3$^{[2]}$ type]\label{prop:k3twosmall}
Assume \textup{(P4)}. Let $X$ be a smooth projective variety of K3$^{[2]}$
type, with $T=T(X)_{\QQ}$ and the form $q$ as in \Cref{prop:k3ntype}, and
suppose that $\NS(X)_{\QQ}$ contains a nonzero class of square $0$ or a
class of square $-2$ for $q$. This holds whenever $\rho(X)\ge4$, that is
$\dim T\le19$, and whenever $X$ admits a Lagrangian fibration. Then the
Hodge conjecture holds for $X$. In particular it holds for $F(Y)$ and for
$Y\times Y$ for every smooth cubic fourfold $Y$ for which $F(Y)$ satisfies
the hypothesis, for instance when $\dim T(Y)\le19$.
\end{proposition}

\begin{proof}
Let $\tilde\Lambda=H^{2}(X,\QQ)\oplus\QQ e$ with $q(e)=2$. By \cite{Bea83},
$H^{2}(X,\ZZ)$ is isometric to $\Lambda\oplus\langle-2\rangle$, where
$\Lambda$ is the K3 lattice, and $\langle-2\rangle\oplus\langle2\rangle$
is a hyperbolic plane over $\QQ$, so $\tilde\Lambda$ is isometric to
$\Lambda_{\QQ}\oplus U_{\QQ}$ with $U_{\QQ}$ a hyperbolic plane. The
orthogonal complement of $T$ in $\tilde\Lambda$ is
$N'=\NS(X)_{\QQ}\oplus\QQ e$. A nonzero vector $\alpha+ae$ of $N'$ with
$q(\alpha)+2a^{2}=0$ is a nonzero isotropic $\alpha$ when $a=0$ and gives
$q(\alpha/a)=-2$ otherwise, so the hypothesis says exactly that $N'$ is
isotropic. An isotropic vector of $N'$ lies in a hyperbolic plane
$H\subseteq N'$ \cite{OMe00}, and by Witt's cancellation theorem the
orthogonal complement of $H$ in $\tilde\Lambda$, which contains $T$, is
isometric to $\Lambda_{\QQ}$. This gives an isometric embedding
$\iota\colon T\to\Lambda_{\QQ}$. When $\rho(X)\ge4$, the space $N'$ has
dimension at least five and signature $(2,\rho(X)-1)$, since $q$ has
signature $(1,\rho(X)-1)$ on $\NS(X)_{\RR}$, which contains an ample class,
so $N'$ is isotropic by the Hasse--Minkowski theorem \cite[\S66]{OMe00}. If
$f\colon X\to B$ is a Lagrangian fibration and $h$ is ample on the surface
$B$, then $f^{*}h\ne0$ and $\int_{X}(f^{*}h)^{4}=0$, so the relation
$\int_{X}\alpha^{4}=3q(\alpha)^{2}$ of Fujiki and Beauville (proof of
\Cref{prop:k3ntype}(iii)) gives $q(f^{*}h)=0$.

The line $\ell=\iota(H^{2,0}(X))$ satisfies
$(\ell,\ell)=0$ and $(\sigma,\bar\sigma)>0$ for $0\ne\sigma\in\ell$, so there
is a K3 surface $S$ with a marking that carries $H^{2,0}(S)$ to $\ell$
\cite[Chs.~6--7]{Huy16}. A rational class orthogonal to $\ell$ has its
component in $\iota(T)$ of type $(1,1)$, and $T$ contains no nonzero rational
class of type $(1,1)$, since such a class would lie in $\NS(X)_{\QQ}$, on
which $q$ is nondegenerate. So $\NS(S)_{\QQ}=\iota(T)^{\perp}$, of signature
$(1,21-\dim T)$, which contains a class of positive square, so $S$ is
projective \cite[Ch.~8]{Huy16}, and $\iota$ is a Hodge isometry from $T$ onto
$T(S)_{\QQ}$.

Let $M=S^{[2]}$. By \cite{Bea83}, $H^{2}(M,\QQ)$ is the orthogonal sum of
$H^{2}(S,\QQ)$, embedded as a sub-Hodge structure on which the
Beauville--Bogomolov form of $M$ is the intersection form, and the line of
half the class of the exceptional divisor; so $T(M)_{\QQ}=T(S)_{\QQ}$ and
$\iota$ is a Hodge isometry from $T(X)_{\QQ}$ onto $T(M)_{\QQ}$. The
varieties $X$ and $M$ are deformation equivalent, so $H^{2}(X,\QQ)$ and
$H^{2}(M,\QQ)$ are isometric, and by Witt's cancellation theorem so are
$\NS(X)_{\QQ}$ and $\NS(M)_{\QQ}$. The sum $f$ of such an isometry and
$\iota$ is a rational Hodge isometry $H^{2}(X,\QQ)\to H^{2}(M,\QQ)$, the
N\'eron--Severi parts being of type $(1,1)$ on both sides. By
\cite[Theorem~1.1]{Mar24}, applied to $f$ and to $f^{-1}$, and the argument
of the proof of \Cref{prop:k3ntype}(ii), with the K\"unneth projectors of
$X$ and $M$ algebraic by \cite[Theorem~1.1]{CM13}, there are algebraic
classes on $X\times M$ and on $M\times X$ that induce nonzero rational
multiples of $f$ and of $f^{-1}$ on $H^{2}$ and zero in the other degrees.

If $E=\End_{\mathrm{Hdg}}(T)$ is $\QQ$ or a CM field,
\Cref{prop:k3ntype}(iii) gives the Hodge conjecture for $X$. Let $E$ be
totally real and $\varphi\in E$. Conjugation by $f$ carries $E$ onto
$\End_{\mathrm{Hdg}}(T(M)_{\QQ})$; let $\varphi_{M}=f\varphi f^{-1}$ on
$T(M)_{\QQ}$. The Hodge conjecture holds for $M$ by \Cref{cor:k3all}(ii), so
by \Cref{prop:k3ntype}(iii) for $M$ the endomorphism $\varphi_{M}\oplus0$ of
$H^{2}(M,\QQ)$ is induced by an algebraic cycle on $M\times M$. Since $f$
carries $\NS(X)_{\QQ}$ to $\NS(M)_{\QQ}$ and $T$ to $T(M)_{\QQ}$,
$f^{-1}\circ(\varphi_{M}\oplus0)\circ f=\varphi\oplus0$, and the composite
of the three correspondences \cite[\S16.1]{Ful98} induces a nonzero multiple
of it. So every element of $E$ acts on $T$ through an algebraic cycle on
$X\times X$, and \Cref{prop:k3ntype}(iii) gives the Hodge conjecture for
$X$.

For a cubic fourfold $Y$, the Fano variety $F(Y)$ is of K3$^{[2]}$ type and
$T(F(Y))_{\QQ}=\alpha(T(Y))$ has the dimension of $T(Y)$, by the proof of
\Cref{cor:fanolines}; so the Hodge conjecture holds for $F(Y)$, and by
\Cref{cor:fanolines} for $Y\times Y$.
\end{proof}

The hypothesis is also necessary for this argument: if $T$ embeds
isometrically in $\Lambda_{\QQ}$, then by Witt's extension theorem its
orthogonal complement in $\tilde\Lambda$ is isometric to that of the image,
which contains $U_{\QQ}$, so $N'$ is isotropic. The varieties that
\Cref{prop:k3twosmall} leaves are therefore those with $\rho(X)\le3$ whose
N\'eron--Severi space is anisotropic and does not represent $-2$; for
$\rho(X)=1$ this is every $X$, since $\NS(X)_{\QQ}=\QQ h$ with $q(h)>0$.
With real multiplication by a field of degree $r\ge2$,
$23-\rho(X)=\dim T=r\dim_{E}T$ with $\dim_{E}T\ge3$, so the shapes left
are $(r,\dim_{E}T)$ among $(2,10)$, $(4,5)$, $(5,4)$ for $\rho(X)=3$,
$(3,7)$, $(7,3)$ for $\rho(X)=2$, and $(2,11)$ for $\rho(X)=1$.

\subsection{Fermat and diagonal varieties}\label{ssec:fermatall}

\begin{corollary}[Hypersurfaces of simplex type]\label{cor:simplexhc}
Assume \textup{(P5)}. The Hodge conjecture holds on every smooth hypersurface
of simplex type in a projective toric variety (\Cref{thm:simplextype}), so on
every smooth Delsarte hypersurface of every dimension and on every cyclic
cover of \Cref{cor:delsarteall}(ii).
\end{corollary}

\begin{proof}
A smooth hypersurface $\sum_{i}a_{i}x_{i}^{e}=0$ has every $a_{i}\ne0$ and
becomes the Fermat hypersurface after rescaling the coordinates, so (P5)
gives the Hodge conjecture on every Fermat variety $X^{j}_{e}$ with $e\ge2$;
for $e=1$ it is a projective space. \Cref{thm:simplextype}(ii) and
\Cref{cor:delsarteall} conclude.
\end{proof}

\begin{corollary}[Very general diagonal complete intersections]
\label{cor:vgalldegrees}
Assume \textup{(P5)}, and keep the notation of \Cref{thm:vandermonde}, with
$d\ge2$, $N\ge2$, $1\le c\le N-1$ and $r=N-c$. For $\lambda\in U$ outside a
countable union of proper Zariski closed subsets, the Hodge conjecture holds
on every power of $X_{d,r}(\lambda)$, so on every power of every complete
intersection of Vandermonde type with these $\lambda_{i}$.
\end{corollary}

\begin{proof}
The curve $C=X_{d,1}(\lambda)$ is the connected abelian cover of $\PP^{1}$
with group $H$, branched over the $\lambda_{i}$, whose monodromy sends a loop
around $\lambda_{i}$ to the image of the $i$-th basis vector, which is not
the identity (\Cref{thm:vandermonde}). As $\lambda$ runs over $U$ these
curves are the members of the full marked family of that cover: every
configuration of $N+1$ points of $\PP^{1}$ is M\"obius equivalent to one in
$\CC$, and the map from $U$ to the configurations up to M\"obius
transformations is a surjective submersion with connected fibres. By
\Cref{thm:cdk} and \Cref{lem:hodgeanalytic}, applied to the variations
$H^{1}(C)^{\otimes2a}$ over $U$, the points $\lambda$ at which some rational
tensor of type $(a,a)$ is not of that type along all of $U$ lie in a
countable union of proper Zariski closed subsets; off it, such a tensor is
of type $(a,a)$ on the whole marked family, so the parameter of $C$ is
Hodge generic and (P5) gives the Hodge conjecture on every power $C^{m}$. By
\Cref{thm:vandermonde}(ii), $X_{d,r}(\lambda)^{u}\cong(C^{r})^{u}/G^{u}$
through $\Phi^{u}$. For a Hodge class $x$ on $X_{d,r}(\lambda)^{u}$, the class
$(\Phi^{u})^{*}x$ is a Hodge class on $C^{ru}$, hence algebraic, and
$x=(\Phi^{u})_{*}(\Phi^{u})^{*}x/|G|^{u}$.
\end{proof}

For $c\le2$, which includes every smooth complete intersection of two
diagonal hypersurfaces of the same degree, \Cref{cor:vgalldegrees} is the
second statement of (P5), proved in \cite{OAI26a} by a Gale duality that
makes the curve appear for two equations; \Cref{thm:vandermonde} supplies
the curve for every $c$. For $d=2$ and $r$ even the Hodge conjecture
on $X_{d,r}(\lambda)$ itself is \Cref{thm:vgvandermonde}(iii), proved in
this paper without the preprints, and so are its cases $d\in\{3,4\}$ with
$r\le6$ and $d=6$ with $r\le4$.

\subsection{CM points and what remains}\label{ssec:remains}

\begin{corollary}[(F2) with every CM point algebraic]\label{cor:f2cm}
Assume \textup{(P1)}.
\begin{enumerate}[label=\textup{(\roman*)},leftmargin=2.6em,itemsep=2pt]
\item For every polarised abelian scheme $\pi\colon\cA\to S$ over a smooth
  connected complex quasi-projective variety and every global section $\xi$
  of $R^{2p}\pi_{*}\QQ$ that is a Hodge class at one point, the algebraic
  locus $\Sigma(\xi)$ contains every point $s$ at which $\cA_{s}$ is of CM
  type.
\item \textup{(F2)} holds if and only if, for every such $\pi$ with a member
  of CM type and every such $\xi$, the locus $\Sigma(\xi)$ has an interior
  point.
\item For every CM field $F$, $n\ge2$ and $\delta$, the algebraic locus
  $\Sigma_{F}$ of \Cref{thm:main} contains the CM points of $\cD_{F}$, a dense
  set contained in the special locus $\Sigma^{\mathrm{sp}}_{F}$, which is
  meagre.
\end{enumerate}
\end{corollary}

\begin{proof}
(i) By the theorem of the fixed part, as in the proof of
\Cref{thm:f2propagation}, $\xi_{s}$ is a Hodge class at every $s$, and it is
algebraic at the points of CM type by (P1). (ii) By (i) the hypothesis of
\Cref{thm:f2propagation}(v) holds for every such $\xi$, so (ii) is the
equivalence of (i) and (v) there. (iii) The CM points are dense by
\Cref{lem:cmdense}, lie in $\Sigma_{F}$ by (i) and in
$\Sigma^{\mathrm{sp}}_{F}$ by \Cref{prop:speciallocus}(iii), which is meagre
by \Cref{prop:speciallocus}(ii).
\end{proof}

So (P1) does not decide the dichotomy of \Cref{thm:main}(iii): the CM points
share every property that \Cref{thm:escape} shows to be insufficient, and an
interior point of $\Sigma_{F}$ still requires the Weil classes at a member
with full Hodge group, which is what (P2) provides for the split eightfolds.

The abelian covers of (P5) reach the Weil classes in a second way: a
character of the covering group cuts out of the Jacobian a piece on which a
cyclotomic field acts, and when that piece is of Weil type and its periods
fill an open subset of $\cD_{F}$, (P5) and \Cref{thm:main}(iii) close the
algebraic locus. A count of moduli shows that this happens only in low
dimension. Write $\varphi$ for Euler's function.

\begin{proposition}[Weil pieces of abelian covers]\label{prop:coversreach}
In the setting of \textup{(P5)}, let $C_{s}\to B_{s}$, $s\in T$, be the full
marked family of a connected abelian cover with group $P$ of a curve of
genus $g$, branched over a set $D$ whose meridians have nontrivial image in
$P$. Let $\chi$ be a character of $P$ of order $\ell\ge3$, put
$F=\QQ(\zeta_{\ell})$, of degree $2m=\varphi(\ell)$, let $D_{\chi}\subseteq D$
be the set of branch points whose meridian does not lie in $\ker\chi$,
$k_{\chi}=|D_{\chi}|$, suppose $2g-2+k_{\chi}>0$, and let
$H_{\chi}\subseteq H^{1}(C_{s},\QQ)$ be the rational subspace whose
complexification is the sum of the eigenspaces of $P$ for the characters
$\chi^{j}$, $j$ prime to $\ell$.
\begin{enumerate}[label=\textup{(\roman*)},leftmargin=2.6em,itemsep=2pt]
\item $H_{\chi}$ is a sub-Hodge structure, $P$ acts on it through an action
  of $F$, and each of the $2m$ eigenspaces has dimension $2g-2+k_{\chi}$.
\item The polarised variation $H_{\chi}$ over $T$ is pulled back from the
  marked family of curves of genus $g$ with $k_{\chi}$ marked points, a
  complex manifold of dimension $3g-3+k_{\chi}$.
\item If $H_{\chi}$ is of $(F,n)$-Weil type, then $2n=2g-2+k_{\chi}$, and the
  image of $T$ in $\cD_{F}$, which has dimension $mn^{2}$, is the image of a
  complex manifold of dimension $3n-k_{\chi}/2$. It is meagre when $m\ge2$
  and $n\ge2$, when $m=1$ and $n\ge4$, and when $m=1$, $n=3$ and
  $k_{\chi}>0$.
\item If $m=1$, $n=3$ and $k_{\chi}=0$, so that $C_{s}/\ker\chi\to B_{s}$ is
  an \'etale cyclic cover of degree $\ell\in\{3,4,6\}$ of a curve of genus
  four given by a line bundle $\eta$ of order $\ell$, and if for one $s$ the
  multiplication map
  $H^{0}(K_{B_{s}}\otimes\eta)\otimes H^{0}(K_{B_{s}}\otimes\eta^{-1})\to
  H^{0}(K_{B_{s}}^{\otimes2})$ is surjective, then \textup{(P5)} gives
  $\Sigma_{F}=\cD_{F}$ for the family of $(F,3)$-Weil type that contains
  $H_{\chi}$.
\end{enumerate}
\end{proposition}

\begin{proof}
(i) The characters $\chi^{j}$ with $j$ prime to $\ell$ form one orbit of
$\Gal(\bar\QQ/\QQ)$, so the sum of their eigenspaces is defined over $\QQ$,
and it is a sub-Hodge structure because $P$ acts by automorphisms of the
curve. A generator $\gamma$ of $P/\ker\chi\cong\mu_{\ell}$ acts on each
eigenspace by a primitive $\ell$-th root of unity, so the cyclotomic
polynomial $\Phi_{\ell}$ kills $\gamma$ on $H_{\chi}$, and $\QQ[\gamma]\cong
F$ acts, compatibly with the polarisation since $P$ preserves the
intersection form. Put $C_{\chi}=C_{s}/\ker\chi$; pull-back identifies the
$\chi^{j}$-eigenspaces of $H^{1}(C_{\chi},\CC)$ and of $H^{1}(C_{s},\CC)$,
since $\chi^{j}$ is trivial on $\ker\chi$. The cyclic cover $C_{\chi}\to
B_{s}$ is branched exactly over $D_{\chi}$. Let $U=B_{s}\setminus D_{\chi}$
and $\tilde U\subseteq C_{\chi}$ its preimage. The direct image of
$\CC_{\tilde U}$ is the sum of rank one local systems $L_{\psi}$ over the
characters $\psi$ of $\mu_{\ell}$, each with Euler characteristic
$2-2g-k_{\chi}$; for $\psi=\chi^{j}$ the local system is nontrivial, so
$H^{0}(U,L_{\psi})=0$, and $H^{2}(U,L_{\psi})$ vanishes, because $U$ is not
compact when $k_{\chi}>0$ and by Poincar\'e duality with the nontrivial dual
system when $k_{\chi}=0$. So $\dim H^{1}(\tilde U,\CC)_{\chi^{j}}=
2g-2+k_{\chi}$. In the exact sequence
$0\to H^{1}(C_{\chi})\to H^{1}(\tilde U)\to\bigoplus_{x}\CC\to H^{2}(C_{\chi})$,
the sum over the points over $D_{\chi}$, the group $\mu_{\ell}$ permutes the
points over a branch point $p$ with stabiliser generated by the image of the
meridian at $p$, on which $\chi^{j}$ is nontrivial, so the
$\chi^{j}$-component of the third term is zero, and
$H^{1}(C_{\chi})_{\chi^{j}}=H^{1}(\tilde U)_{\chi^{j}}$.

(ii) The meridians at the points of $D\setminus D_{\chi}$ lie in $\ker\chi$,
so the homomorphism $\pi_{1}(B_{s}\setminus D)\to\mu_{\ell}$ factors through
$\pi_{1}(U)$ and $C_{\chi}$ is the cover of $B_{s}$ determined by the marked
curve $(B_{s},D_{\chi})$ and that topological datum. Hence the Hodge
structure on $H^{1}(C_{\chi})$, its polarisation and the action of $\mu_{\ell}$
depend only on the image of $s$ in the marked family of curves of genus $g$
with $k_{\chi}$ points, which has dimension $3g-3+k_{\chi}$ when
$2g-2+k_{\chi}>0$; and the polarisation of $H_{\chi}\subseteq H^{1}(C_{s})$
is $|\ker\chi|$ times the one pulled back from $C_{\chi}$.

(iii) The space $H_{\chi}$ has dimension $2m(2g-2+k_{\chi})$ over $\QQ$, so
$2n=2g-2+k_{\chi}$ over $F$, and
$3g-3+k_{\chi}=3(n-k_{\chi}/2)+k_{\chi}=3n-k_{\chi}/2$. The marking makes the
period map of $H_{\chi}$ a holomorphic map to $\cD_{F}$, which is a product
of $m$ hermitian symmetric domains of $\mathrm{U}(n,n)$, each of dimension
$n^{2}$. The image of a holomorphic map from a manifold of smaller dimension
is a countable union of compact sets of measure zero, so it is meagre. Now
$mn^{2}\ge2n^{2}>3n$ for $m\ge2$ and $n\ge2$, $n^{2}>3n$ for $n\ge4$, and
$9>9-k_{\chi}/2$ for $n=3$ and $k_{\chi}>0$.

(iv) Here $H^{1,0}$ of the $\chi$-eigenspace is $H^{0}(B_{s},K_{B_{s}}\otimes
\eta)$ for a suitable choice of $\eta$ among $\eta^{\pm1}$, of dimension
three, and the tangent space of $\cD_{F}$ is
$\Hom\bigl(H^{0}(K\otimes\eta),H^{1}(\eta)\bigr)$, of dimension nine. By
Griffiths' theorem on the derivative of the period map \cite[Chapter~10]{Voi02},
applied to the $\chi$-eigenspace of $C_{\chi}$, whose Kodaira--Spencer class
is pulled back from $B_{s}$, the differential at $s$ sends a class in
$H^{1}(B_{s},T_{B_{s}})$ to cup product with it, and its
transpose, after Serre duality $H^{1}(\eta)\cong H^{0}(K\otimes\eta^{-1})^{\vee}$,
is the multiplication map of the statement. Both sides have dimension nine,
so when that map is surjective the period map is a local isomorphism at $s$
and its image contains an open set $\Omega\subseteq\cD_{F}$. By (P5) the
Weil classes are algebraic at the image of every parameter outside a
countable union of proper closed analytic subsets of $T$, and the image of
that complement is not meagre in $\Omega$. So $\Sigma_{F}$ is not meagre, and
$\Sigma_{F}=\cD_{F}$ by \Cref{thm:main}(iii).
\end{proof}

Within the reach of (P5), \Cref{prop:coversreach} leaves only the Weil
fourfolds, where Markman's theorem already applies, and the sixfolds of
part (iv), which (P2) covers by \Cref{thm:weilsix}. The eightfolds of
nontrivial discriminant, which \Cref{prop:stabilise} reduces to the split
tenfolds, and the Weil classes of the fields of degree at least four, are not
reached by any piece $H_{\chi}$ except on a meagre set.

\begin{remark}[What remains of (F2) and (F3$'$)]\label{rem:remains}
With (P1) to (P5) and the results above, the open cases named in
\Cref{thm:twostand} change as follows. In part (ii) the non-split Weil
sixfolds and the squares of Mumford fourfolds are settled
(\Cref{thm:weilsix}, \Cref{cor:mumfordall}); in part (iii) the locus of the
real multiplication is the whole domain in every family of \Cref{setup:rm}
(\Cref{cor:k3all}). Of the three frontier cases of \textup{(F3$'$)} in
\Cref{rem:f3primefrontier}, the first and the third are settled
(\Cref{cor:k3all}); the second, the fourfolds of K3$^{[2]}$ type with real
multiplication, among them $F(Y)$ and $Y\times Y$ for a cubic fourfold $Y$
with real multiplication (\Cref{cor:fanolines}), is settled when the
N\'eron--Severi space is isotropic or represents $-2$, in particular when
$\rho\ge4$ or there is a Lagrangian fibration (\Cref{prop:k3twosmall}),
and is open for $\rho\le3$ otherwise, in the six shapes listed after that
proposition.

By \Cref{prop:f2isab}, (F2) is the Hodge conjecture for all abelian
varieties. It is now known, or conditional only on (P1) to (P5), for the
abelian varieties of CM type and their powers (P1), for those of dimension at
most five \cite{Mar25a,Mar25b,MZ99}, for the Weil classes of every abelian
sixfold of Weil type and of every split eightfold of an imaginary quadratic
field (\Cref{thm:weilsix}), for the powers of the abelian varieties of
orthogonal type of \Cref{cor:kstype}, among them the Mumford fourfolds
(\Cref{cor:mumfordall}), and for the powers of the Jacobians of Hodge
generic abelian covers (P5). The first cases outside all of these are the
Weil classes of an imaginary quadratic field on abelian eightfolds of
nontrivial discriminant, which \Cref{prop:stabilise} reduces to the split
tenfolds, and on abelian $2n$-folds with $n\ge5$; the Weil classes of a CM
field of degree at least four at the members that are not of CM type, the
first being the eightfolds of a quartic field (\Cref{thm:quarticrank}); and
the exceptional Hodge classes of abelian varieties whose Hodge group is not
of any of these types. \textup{(F3$'$)} is open beyond $\mathcal{A}$ and the
cases of \Cref{rem:f3primefrontier}, \Cref{cor:k3all} and
\Cref{prop:k3twosmall}, and the step of
\Cref{rem:f3primeopen}, a new correspondence for a variety that no known
correspondence links to an abelian variety, is not touched by any of these
results. By \Cref{thm:final}(vi) the Hodge conjecture is equivalent to (F2)
and \textup{(F3$'$)} together, and neither is proved, with or without the
preprints.
\end{remark}
